% =============================================================================
% Referee-revised patch v2: §10 Aggregation and zero-free strip + §11 RH
%   sec:aggregation-and-zero-free-strip
%   sec:riemann-hypothesis
%
% Action: replace current §§10--11 with the bodies below; in the same
%         merge, update §1.2 thm:master-conditional to list the new
%         Hypothesis hyp:finite-height-verification among its inputs
%         (adjacent edit; see "Master-conditional adjustment" block at
%         the end of this file).
%
% Referee hardening relative to the submitted patch:
%   - fixes the location of branch exhaustion (it lives in §8, not §6);
%   - adds an explicit high-height branch-entry input, since branch
%     exhaustion only applies to admissible packets after branch entry;
%   - separates the final single-zero assembly from any
%     polynomial-weight or N-point scalar aggregation needed to prove
%     the rigidity/suppression inputs;
%   - makes the finite-height step explicit, because an ineffective
%     T_0 requires a finite verification certificate up to that same
%     T_0;
%   - records the two routes to remove the finite-height input: prove
%     all-height applicability of the local argument, or make T_0
%     effective and place it below a certified numerical verification
%     height (Platt--Trudgian 3*10^12).
%
% Status (§10):
%   \statusmig      not-started -> migrated
%   \statusreferee  not-started -> in-progress
%   \statussympy    not-started -> n/a
%   \statussim      not-started -> n/a
%   \statuslean     not-started -> n/a
%
% Status (§11):
%   \statusmig      n/a
%   \statusreferee  not-started -> in-progress
%   \statussympy    n/a
%   \statussim      n/a
%   \statuslean     not-started
%
% Closes:
%   rem:wip-zero-free-strip
% =============================================================================

% =============================================================================
\section{Aggregation and zero-free strip}
\label{sec:aggregation-and-zero-free-strip}
\statuspaper{LIVE}\quad
\statusmig{migrated}\quad
\statusreferee{in-progress}\quad
\statussympy{n/a}\quad
\statussim{n/a}\quad
\statuslean{n/a}
\par\medskip

The preceding sections separate local exclusion from global assembly.  The
local comparison theorems of Section~\ref{sec:projective-rigidity} exclude
admissible packets in the two-point scalar branch or in the four-point
full-quartet branch.  The branch-exhaustion input
(Hypothesis~\ref{hyp:branch-exhaustion}) says that these two branches cover
all admissible off-line packets produced by the branch-entry construction.
The present section records the final high-height assembly.  It does not
prove branch entry, actual-zeta suppression, source-energy transfer, or the
analytic estimates needed to consume the \(N\)-point cancellation of
Section~\ref{sec:n-point-odd-moment-cancellation}; those inputs must already
be included in the rigidity/suppression theorems being invoked.

\begin{hypothesis}[High-height branch entry]
\label{hyp:high-height-branch-entry}
There exists a height \(T_{\mathrm{entry}}>0\) such that every hypothetical
off-line zero
\[
        \rho=\beta+i\gamma,
        \qquad \beta\ne \tfrac12,
        \qquad |\gamma|\ge T_{\mathrm{entry}},
\]
produces an admissible off-line packet in the local jet--Gram framework, with
all normalizations, retained subdomains, and loss ledgers required by
Hypothesis~\ref{hyp:branch-exhaustion} and by the applicable rigidity theorem.
\end{hypothesis}

\begin{remark}[All-height alternative versus finite verification]
\label{rem:all-height-vs-finite-verification}
The high-height formulation is a proof-management device, not an intrinsic
feature of the desired theorem.  There are two possible ways to remove the
finite-height input in Section~\ref{sec:riemann-hypothesis}.

First, one may strengthen Hypothesis~\ref{hyp:high-height-branch-entry} and
the local rigidity inputs so that they apply at every nontrivial-zero height,
with no asymptotic lower-height cutoff.  This would require replacing the
large-\(T\) estimates used in the local chart, whitening, toy high-frequency
limit, and packet-retention arguments by uniform finite-height estimates.  In
that case one may take \(T_{\mathrm{entry}}=0\) and no numerical finite-height
verification is needed for the nontrivial zeros.

Second, one may keep the high-height proof but make the resulting threshold
\(T_0\) effective.  If a certified computation verifies that all zeros with
\(0<|\gamma|\le H_{\mathrm{cert}}\) lie on the critical line and if
\(T_0\le H_{\mathrm{cert}}\), then the finite-height verification input of
Hypothesis~\ref{hyp:finite-height-verification} follows from that certificate.
For example, Platt--Trudgian verify the Riemann hypothesis rigorously up to
height \(3\cdot10^{12}\) by interval arithmetic.  If \(T_0\) is ineffective,
or if \(T_0>H_{\mathrm{cert}}\), then a separate finite-height certificate up
to \(T_0\) is still required.
\end{remark}

\begin{remark}[Role of the \(N\)-point section in the final assembly]
\label{rem:n-point-role-in-aggregation}
Section~\ref{sec:n-point-odd-moment-cancellation} supplies algebraic scalar
moment cancellation.  If a two-point or four-point suppression theorem uses
that cancellation, then the corresponding analytic remainder estimates,
polynomial-weight density losses, and exponent bookkeeping must be absorbed in
that suppression or rigidity theorem before Theorem~\ref{thm:zero-free-strip}
is invoked.  The final assembly below does not itself turn the
\(N\)-point algebra into an analytic suppression estimate.
\end{remark}

\begin{theorem}[Zero-free strip above an absolute height]
\label{thm:zero-free-strip}
Assume Hypothesis~\ref{hyp:high-height-branch-entry},
Hypothesis~\ref{hyp:branch-exhaustion}, and the inputs of
Theorems~\ref{thm:two-point-rigidity} and~\ref{thm:four-point-rigidity}
with absolute high-height thresholds.  Then there exists an absolute constant
\(T_0>0\) such that
\begin{equation}
\label{eq:zero-free-strip}
        \zeta(\beta+i\gamma)\ne 0
        \qquad
        \text{for every } \beta\ne\tfrac12 \text{ and } |\gamma|\ge T_0.
\end{equation}
\end{theorem}

\begin{proof}
Let \(T_0^{(2)}\) and \(T_0^{(4)}\) be the high-height thresholds for the
two-point and four-point rigidity theorems, respectively, and let
\[
        T_0:=\max\{T_{\mathrm{entry}},T_0^{(2)},T_0^{(4)}\}.
\]
Suppose, for contradiction, that
\(\rho=\beta+i\gamma\) is a nontrivial zero of \(\zeta\) with
\(\beta\ne\tfrac12\) and \(|\gamma|\ge T_0\).  By
Hypothesis~\ref{hyp:high-height-branch-entry}, \(\rho\) produces an
admissible off-line packet.  By Hypothesis~\ref{hyp:branch-exhaustion}, this
packet is excluded by either the two-point scalar branch of
Theorem~\ref{thm:two-point-rigidity} or the four-point full-quartet branch of
Theorem~\ref{thm:four-point-rigidity}.

If the two-point branch applies, then \(|\gamma|\ge T_0^{(2)}\), and
Theorem~\ref{thm:two-point-rigidity}, under its stated inputs, contradicts the
existence of \(\rho\).  If the four-point branch applies, then
\(|\gamma|\ge T_0^{(4)}\), and Theorem~\ref{thm:four-point-rigidity}, under
its stated inputs, contradicts the existence of the symmetry quartet generated
by \(\rho\).  Both alternatives are impossible.  Hence
\eqref{eq:zero-free-strip} holds.
\end{proof}

\begin{remark}[No additional zero-counting loss]
\label{rem:zero-free-strip-no-aggregation-loss}
Once Hypothesis~\ref{hyp:high-height-branch-entry},
Hypothesis~\ref{hyp:branch-exhaustion}, and the relevant local rigidity theorem
are available with absolute thresholds, the final high-height assembly does not
accumulate a loss over the set of zeros.  Each hypothetical off-line zero above
\(T_0\) is contradicted individually.  Any polynomial-weight density loss,
\(N\)-point moment-cancellation loss, or source-energy loss must already have
been charged in the proof of the suppression or rigidity input being used.
\end{remark}

\begin{remark}[Scope of \S\ref{sec:aggregation-and-zero-free-strip}]
\label{rem:zero-free-strip-scope}
This section proves only the formal high-height implication from branch entry,
branch exhaustion, and local rigidity to a zero-free strip.  It does not address
heights \(|\gamma|<T_0\), and it does not make \(T_0\) effective.  The finite
height range is a separate input in Section~\ref{sec:riemann-hypothesis} unless
all constants have been made explicit and lie below an already certified
verification height.
\end{remark}

% =============================================================================
\section{The Riemann hypothesis}
\label{sec:riemann-hypothesis}
\statuspaper{LIVE}\quad
\statusmig{n/a}\quad
\statusreferee{in-progress}\quad
\statussympy{n/a}\quad
\statussim{n/a}\quad
\statuslean{not-started}
\par\medskip

Theorem~\ref{thm:zero-free-strip} excludes off-line zeros above the height
\(T_0\) supplied by the high-height assembly.  To conclude the full Riemann
hypothesis, the bounded interval \(|\gamma|<T_0\) must also be certified.  If
\(T_0\) is not effective, this finite-height certification remains a separate
input.

\begin{hypothesis}[Finite-height verification]
\label{hyp:finite-height-verification}
For the height \(T_0\) appearing in Theorem~\ref{thm:zero-free-strip}, every
nontrivial zero \(\rho=\beta+i\gamma\) of \(\zeta\) with \(|\gamma|<T_0\)
satisfies \(\beta=\tfrac12\).
\end{hypothesis}

\begin{theorem}[Conditional Riemann hypothesis]
\label{thm:riemann-hypothesis}
Assume the inputs of Theorem~\ref{thm:zero-free-strip} and
Hypothesis~\ref{hyp:finite-height-verification}.  Then every nontrivial zero
of \(\zeta\) satisfies \(\RePart s = \tfrac12\).
\end{theorem}

\begin{proof}
Let \(\rho=\beta+i\gamma\) be a nontrivial zero of \(\zeta\).  If
\(|\gamma|\ge T_0\), then \(\beta=\tfrac12\) by
Theorem~\ref{thm:zero-free-strip}.  If \(|\gamma|<T_0\), then
\(\beta=\tfrac12\) by Hypothesis~\ref{hyp:finite-height-verification}.
Thus every nontrivial zero lies on the critical line.
\end{proof}

\begin{remark}[Effectivity and numerical verification]
\label{rem:finite-height-effectivity}
Certified finite-height verification may be used once the height threshold is
made explicit.  Let \(H_{\mathrm{cert}}\) denote a rigorous certified height up
to which every nontrivial zero has been verified to lie on the critical line.
If the proof produces an explicit \(T_0\le H_{\mathrm{cert}}\), then
Hypothesis~\ref{hyp:finite-height-verification} follows from that certificate.
Known computations include the rigorous interval-arithmetic verification of
Platt--Trudgian up to \(H_{\mathrm{cert}}=3\cdot10^{12}\).  If the proof only
produces a non-effective \(T_0\), then known finite computations cannot be
invoked automatically; one must either make \(T_0\) effective and compare it
with a certified height, prove all-height applicability of the local argument,
or supply a finite verification certificate for the particular \(T_0\).
\end{remark}

% =============================================================================
% Master-conditional adjustment (adjacent edit, applied at merge time)
%
% Target: rh/rh_rebuild.tex, §1.2 (subsec:master-conditional),
%         theorem block thm:master-conditional.
%
% Reason: §11 thm:riemann-hypothesis is now Conditional, with
%         Hypothesis hyp:finite-height-verification as an explicit input.
%         For consistency, §1.2 thm:master-conditional must list the
%         same hypothesis among its inputs, since it concludes the same
%         RH statement.  Without this adjustment, §1.2 would overclaim
%         relative to §11.
%
% Replace the current input list inside thm:master-conditional with the
% block below (the only change is the new "the finite-height
% verification (Hypothesis~\ref{hyp:finite-height-verification})" item,
% placed after the aggregation step).  No other text in §1.2 changes.
%
% --- old (current rh_rebuild.tex, line ~287--301) ---
%   ... aggregation step (Theorem~\ref{thm:zero-free-strip}).  Then every
%   nontrivial zero of \(\zeta\) satisfies \(\RePart s = \tfrac12\).
%
% --- new ---
%   ... aggregation step (Theorem~\ref{thm:zero-free-strip}); and the
%   finite-height verification
%   (Hypothesis~\ref{hyp:finite-height-verification}).  Then every
%   nontrivial zero of \(\zeta\) satisfies \(\RePart s = \tfrac12\).
% =============================================================================
